% ==============================================================================
% Preamble: [YOUR PAPER TITLE]
% ==============================================================================

\documentclass[12pt,letterpaper]{article}

% --- Encoding & Fonts (XeLaTeX) ---
\usepackage{fontspec}
\setmainfont{Times New Roman}

% --- Page Layout ---
\usepackage[margin=1in]{geometry}
\usepackage{setspace}
% House style is 1.5x body spacing. Use \setstretch{1.5} explicitly rather than
% \onehalfspacing. The latter applies setspace's font-size-dependent stretch
% (about 1.24 at 12pt), which looks closer to 1.25x than to 1.5x.
\setstretch{1.5}

% --- Math ---
\usepackage{amsmath}
\usepackage{amssymb}
\usepackage{amsthm}
\usepackage{bm}

% --- Tables & Figures ---
\usepackage{graphicx}
\usepackage{booktabs}
\usepackage{tabularx}
\usepackage{float}
\usepackage{multirow}
\usepackage{array}
\usepackage{caption}
\captionsetup{font=small,labelfont=bf,format=hang}

% --- Citations ---
\usepackage[authoryear,round]{natbib}
\bibliographystyle{aer}

% --- Optional: separate appendix bibliography ---
% Uncomment the three lines below to give the appendices their own reference
% list. Appendix citations then use \citetapp and \citepapp, the list is printed
% with \bibliographyapp{...}, and the compile pipeline gains one step, bibtex app,
% after bibtex main.
% \usepackage{multibib}
% \newcites{app}{References (Appendices)}
% \bibliographystyleapp{aer}

% --- Colors & Links ---
\usepackage{xcolor}
\usepackage{hyperref}
\hypersetup{
  colorlinks=true,
  linkcolor=blue!60!black,
  citecolor=blue!60!black,
  urlcolor=blue!60!black
}

% --- Misc ---
\usepackage{enumitem}
\usepackage{appendix}
\usepackage[hang,flushmargin]{footmisc}

% ==============================================================================
% Custom Commands: Project-Specific Notation
% ==============================================================================
% Define your project-specific notation here. Examples:
%
% \newcommand{\treat}{T}        % treatment variable
% \newcommand{\outcome}{Y}      % outcome variable
% \newcommand{\effect}{\beta}   % treatment effect
%
% Keeping notation in the preamble ensures consistency across all sections.

% Expectations and operators (universally useful)
\newcommand{\E}{\mathbb{E}}
\newcommand{\Var}{\text{Var}}

% ==============================================================================
% Theorem-like environments
% ==============================================================================
\newtheorem{hypothesis}{Hypothesis}
\newtheorem{proposition}{Proposition}
\newtheorem{prediction}{Prediction}
